%\vspace{0.1cm}
\textbf{Resolução:}

\textbf{Passo 1: Identificar as dimensões de cada variável no sistema $MLT$ (Massa, Comprimento, Tempo).}
\begin{align*}
\text{Vazão volumétrica } (Q): \quad &[Q] = \frac{\text{Volume}}{\text{Tempo}} = L^3 \cdot T^{-1} \\
\text{Queda de pressão/comprimento } \left(\frac{\Delta p}{l}\right): \quad &\left[\frac{\Delta p}{l}\right] = \frac{\text{Força/Área}}{\text{Comprimento}} \\
&= \frac{M \cdot L \cdot T^{-2} / L^2}{L} \\
&= M \cdot L^{-2} \cdot T^{-2} \\
\text{Diâmetro do tubo } (d): \quad &[d] = L \\
\text{Viscosidade do fluido } (\mu): \quad &[\mu] = M \cdot L^{-1} \cdot T^{-1}
\end{align*}

\textbf{Passo 2: Configurar a equação pelo método básico.} \\
Assumimos que a vazão $Q$ se relaciona com as outras variáveis através de um produto de potências, multiplicado por uma constante adimensional $k$:
\[ Q = k \cdot \left(\frac{\Delta p}{l}\right)^\alpha \cdot d^\beta \cdot \mu^\gamma \]

\textbf{Passo 3: Substituir as dimensões na equação.}
\begin{align*}
[Q] &= \left[\frac{\Delta p}{l}\right]^\alpha \cdot [d]^\beta \cdot [\mu]^\gamma \\
L^3 \cdot T^{-1} &= (M \cdot L^{-2} \cdot T^{-2})^\alpha \cdot (L)^\beta \cdot (M \cdot L^{-1} \cdot T^{-1})^\gamma
\end{align*}

Agrupando os termos das dimensões do lado direito:
\[ M^0 \cdot L^3 \cdot T^{-1} = M^{\alpha+\gamma} \cdot L^{-2\alpha+\beta-\gamma} \cdot T^{-2\alpha-\gamma} \]

\textbf{Passo 4: Montar e resolver o sistema de equações.} \\
Igualando os expoentes de $M$, $L$ e $T$ de ambos os lados:
\[
\begin{cases}
\text{M: } 0 = \alpha + \gamma \implies \gamma = -\alpha \\
\text{T: } -1 = -2\alpha - \gamma \\
\text{L: } 3 = -2\alpha + \beta - \gamma
\end{cases}
\]

Substituindo $\gamma = -\alpha$ na equação do Tempo (T):
\begin{align*}
-1 &= -2\alpha - (-\alpha) \\
-1 &= -\alpha \implies \alpha = 1
\end{align*}
Como $\gamma = -\alpha$, então $\gamma = -1$.

Substituindo $\alpha = 1$ e $\gamma = -1$ na equação do Comprimento (L) para encontrar $\beta$:
\begin{align*}
3 &= -2(1) + \beta - (-1) \\
3 &= -2 + \beta + 1 \\
3 &= -1 + \beta \implies \beta = 4
\end{align*}

\textbf{Passo 5: Formular a relação final.} \\
Substituindo os valores encontrados ($\alpha = 1$, $\beta = 4$, $\gamma = -1$) na equação assumida no Passo 2:
\begin{align*}
Q &= k \cdot \left(\frac{\Delta p}{l}\right)^1 \cdot d^4 \cdot \mu^{-1} \\
Q &= k \cdot \frac{d^4}{\mu} \left(\frac{\Delta p}{l}\right)
\end{align*}

Onde $k$ representa a constante. Portanto, demonstramos pela análise dimensional que:
\[ Q = (\text{constante})\left( \frac{d^4}{\mu} \right)\left( \frac{\Delta p}{l} \right) \]
